\documentclass[10pt]{article}
\usepackage[T1]{fontenc}
\usepackage[utf8]{inputenc}
\usepackage[french]{babel}
\usepackage{amsmath,amssymb,mathrsfs,mathpazo}
\usepackage{geometry}
\geometry{a4paper,margin=2cm}
\usepackage{graphicx,array,tabularx,tabularray,colortbl,multirow,multicol,booktabs,enumitem,pifont,ulem,fancybox,diagbox,cancel}
\usepackage{pstricks,pstricks-add,pst-plot,pst-tree,pst-node,pst-text,pst-func,pst-math,pst-bspline,pst-3dplot,pst-eucl,pst-solides3d}
\let\pscalculate\undefined
\usepackage{pst-bezier}
\usepackage{tikz,pgfplots,tkz-tab}
\usetikzlibrary{arrows,arrows.meta,patterns,calc,positioning,trees,decorations.markings,shapes.geometric,shapes.misc,shapes.arrows,intersections,decorations.pathmorphing,decorations.pathreplacing,backgrounds,shadows,babel}
\pgfplotsset{compat=1.18}
\usepgfplotslibrary{fillbetween}
\usepackage[np]{numprint}
\usepackage{listings,makecell,pas-tableur,twoopt}
\let\barre\undefined
\usepackage[pstricks]{tablvar}
\definecolor{darkgreen}{RGB}{0,100,0}
\providecommand{\red}{\color{red}}
\providecommand{\blue}{\color{blue}}
\providecommand{\green}{\color{green}}
\providecommand{\black}{\color{black}}
\providecommand{\white}{\color{white}}
\providecommand{\gray}{\color{gray}}
\providecommand{\cyan}{\color{cyan}}
\providecommand{\magenta}{\color{magenta}}
\providecommand{\yellow}{\color{yellow}}
\providecommand{\violet}{\color{violet}}
\providecommand{\orange}{\color{orange}}
\providecommand{\euro}{\texteuro}
\providecommand{\EUR}{\texteuro}
\providecommand{\ieme}{\textsuperscript{e}}
\providecommand{\degres}{\ensuremath{^\circ}}
\providecommand{\degree}{\ensuremath{^\circ}}
\providecommand{\vv}[1]{\overrightarrow{#1}}
\providecommand{\Coord}[1]{#1}
\providecommand{\decosix}{}
\providecommand{\starredbullet}{\ensuremath{\bullet}}
\pagestyle{empty}
\setlength{\parindent}{0pt}
\usepackage{adjustbox}\DeclareUnicodeCharacter{2236}{:}\DeclareUnicodeCharacter{2190}{\ensuremath{\leftarrow}}
\providecommand{\R}{\mathbb{R}}
\providecommand{\N}{\mathbb{N}}
\providecommand{\Z}{\mathbb{Z}}
\providecommand{\Q}{\mathbb{Q}}
\providecommand{\D}{\mathbb{D}}
\providecommand{\C}{\mathbb{C}}
\providecommand{\e}{{\mathrm{e}}}
\providecommand{\vect}[1]{\overrightarrow{#1}}
\providecommand{\vectt}[1]{\overrightarrow{\mathrm{#1}}}
\providecommand{\vvt}[1]{\overrightarrow{\mathrm{#1}}}
\providecommand{\np}[2][]{{#2~#1}}
\providecommand{\vv}[1]{{\overrightarrow{#1}}}
\renewcommand{\d}{{\mathrm{d}}}
\providecommand{\degres}{{^\circ}}
\providecommand{\degre}{{^\circ}}
\providecommand{\blue}{\color{blue}}
\providecommand{\red}{\color{red}}
\providecommand{\enskip}{\quad}
\providecommand{\llbracket}{\lbrack\!\lbrack}
\providecommand{\rrbracket}{\rbrack\!\rbrack}
\providecommand{\euro}{{\text{€}}}
\providecommand{\no}{{\text{n°}}}
\providecommand{\interval}[2]{{[#1;#2]}}
\providecommand{\intervFO}[2]{{[#1;#2[}}
\providecommand{\dotfill}{\ldots}

\providecommand{\ds}{\displaystyle}
\providecommand{\textding}[1]{#1}
\providecommand{\Oij}{\(\mathrm{O};\vec{\imath},\vec{\jmath}\)}
\providecommand{\Oijk}{\(\mathrm{O};\vec{\imath},\vec{\jmath},\vec{k}\)}
\providecommand{\Ouv}{\(\mathrm{O};\vec{u},\vec{v}\)}
\providecommand{\intervFF}[2]{\ensuremath{[#1;#2]}}
\renewcommand{\euro}{\ifmmode\text{\texteuro}\else\texteuro\fi}
\DeclareUnicodeCharacter{00B0}{\ifmmode^{\circ}\else\textdegree\fi}
\DeclareUnicodeCharacter{20AC}{\ifmmode\text{\texteuro}\else\texteuro\fi}
% Énoncés : Éducation nationale. Transcription et corrigés : APMEP, Denis Vergès et contributeurs. Découpage et classement : CoopMaths, Sébastien Lozano, Éric Elter et contributeurs. Adaptation de la présentation : Maths Explorer.
% CC BY-SA : https://coopmaths.fr/www/about/
% https://coopmaths.fr/alea/static/bac/2022/tex/bac_2022_05_sujet2_metropole_3.tex
\begin{document}
\noindent\textit{APMEP / CoopMaths. CC BY-SA. Présentation adaptée par Maths Explorer.}\par\medskip
\medskip

\begin{minipage}{0.6\linewidth}
On considère un cube ABCDEFGH
et on appelle K le milieu du segment [BC].

On se place dans le repère 
$\left(\text{A}~;~\vect{\text{AB}},~\vect{\text{AD}},~\vect{\text{AE}}\right)$ 
et on considère  le tétraèdre EFGK.

On rappelle que le volume d'un tétraèdre est donné par: 

\[V = \dfrac13 \times \mathcal{B} \times h\]

où $\mathcal{B}$ désigne l'aire d'une base et $h$ la hauteur relative à cette base.
\end{minipage} \hfill
\begin{minipage}{0.35\linewidth}
\psset{unit=1cm}
\begin{pspicture}(5.5,5.8)
\psframe(0.2,0.2)(3.7,3.7)
\psline(3.7,0.2)(5,1.9)(5,5.4)(3.7,3.7)(4.35,1.05)
\psline(5,5.4)(1.5,5.4)(0.2,3.7)
\psline[linestyle=dashed](0.2,0.2)(1.5,1.9)(5,1.9)
\psline[linestyle=dashed](1.5,1.9)(1.5,5.4)
\pspolygon[linestyle=dotted,linewidth=1.25pt](0.2,3.7)(5,5.4)(4.35,1.05)
\uput[dl](0.2,0.2){\small A}\uput[dr](3.7,0.2){\small B}\uput[r](5,1.9){\small C}
\uput[dr](1.5,1.9){\small D}\uput[l](0.2,3.7){\small E}\uput[r](3.7,3.7){\small F}
\uput[ur](5,5.4){\small G}\uput[ul](1.5,5.4){\small H}\uput[dr](4.35,1.05){\small K}
\end{pspicture}
\end{minipage}

\medskip

\begin{enumerate}

\item Préciser les coordonnées des points E, F{}, G et K.
\item Montrer que le vecteur $\vect{n}\begin{pmatrix}\phantom{-}2\\-2\\\phantom{-}1\end{pmatrix}$ est orthogonal au plan (EGK).
\item Démontrer que le plan (EGK) admet pour équation cartésienne : $2x - 2y + z - 1 = 0.$
\item Déterminer une représentation paramétrique de la droite $(d)$ orthogonale au plan (EGK)
passant par F{}.
\item Montrer que le projeté orthogonal L de F sur le plan (EGK) a pour coordonnées $\left(\frac59~;~\frac49~;~\frac79\right)$.
\item Justifier que la longueur LF est égale à $\dfrac23$.
\item Calculer l'aire du triangle EFG. En déduire que le volume du tétraèdre EFGK est égal à $\dfrac16$.
\item Déduire des questions précédentes l'aire du triangle EGK.
\item On considère les points P milieu du segment [EG], M milieu du segment [EK] et N milieu du segment[GK]. Déterminer le volume du tétraèdre FPMN.
\end{enumerate}

\bigskip

\end{document}
